\documentclass[spanish,letterpaper,oneside]{article}

% Datos del documento
\def\documenttitle {Sistemas de ecuaciones, perímetros y robótica}
\def\documentsubtitle {Temas selectos de matemáticas I}
\def\documentsubject {Actividad 11 · Ficha y propuesta de calibración}
\def\documentauthor {Mtro. Martín Jonathan de la Cruz Muñoz}
\def\coursename {Temas selectos de matemáticas I}
\def\coursecode {PFM02}

% Identidad institucional
\def\universityname {Instituto de Investigación, Innovación y Estudios de Posgrado para la Educación del Estado de Nuevo León}
\def\universityfaculty {Maestría en Aprendizaje y Enseñanza de las Matemáticas}
\def\universitydepartment {Temas selectos de matemáticas I}
\def\universitydepartmentimage {img/iiiepe/marca-iiiepe.png}
\def\universitydepartmentimagecfg {height=1.75cm}
\def\universitylocation {Monterrey, Nuevo León}
\def\studentfullname {Mtro. Martín Jonathan de la Cruz Muñoz}
\def\studentid {00874}
\def\studentemail {martin.delacruz@campus.iiiepe.edu.mx}
\def\programperiod {1er tetramestre, marzo--julio de 2026}
\def\programmodality {Modalidad presencial}
\def\programshort {MAEM}

\def\authortable {
  \begin{tabular}{ll}
    \textbf{Alumno:} & \studentfullname \\
    Matrícula: & \studentid \\
    Correo institucional: & \studentemail \\
    Unidad: & \coursename \\
    Programa: & \programshort \\
    Periodo: & \programperiod \\
    & \\
    \multicolumn{2}{l}{Fecha de realización: 13 de octubre de 2026} \\
    \multicolumn{2}{l}{\universitylocation}
  \end{tabular}
}

\input{template}

\usepackage{helvet}
\renewcommand{\familydefault}{\sfdefault}
\usepackage{array}
\usepackage{longtable}
\usepackage{pdflscape}
\usepackage{tikz}
\usepackage[most]{tcolorbox}
\usetikzlibrary{arrows.meta,positioning,calc,fit,backgrounds}

\definecolor{iiiepeNavy}{HTML}{163B5C}
\definecolor{iiiepeBlue}{HTML}{1F6F9F}
\definecolor{iiiepeAqua}{HTML}{20B7A8}
\definecolor{iiiepeOrange}{HTML}{F28C28}
\definecolor{iiiepePaper}{HTML}{F6F8F7}
\definecolor{iiiepeInk}{HTML}{1E2A32}
\def\portraittitlecolor {iiiepeNavy}
\setcitestyle{authoryear,open={(},close={)}}

\newcommand{\insertcoverwatermark}{%
  \AddToShipoutPictureBG*{%
    \begin{tikzpicture}[remember picture,overlay]
      \node[opacity=.09,inner sep=0pt] at (current page.center)
      {\includegraphics[width=.58\paperwidth]{img/iiiepe/logo-puntos.png}};
    \end{tikzpicture}%
  }%
}

\newtcolorbox{reportelecturabox}[1][]{%
  enhanced,breakable,
  colback=iiiepePaper,colframe=iiiepeNavy,
  boxrule=1pt,leftrule=4pt,arc=2pt,
  left=8pt,right=8pt,top=6pt,bottom=6pt,
  before skip=6pt,after skip=10pt,
  fonttitle=\bfseries\color{white},coltitle=white,
  colbacktitle=iiiepeNavy,
  attach boxed title to top left={xshift=8pt,yshift=-2pt},
  boxed title style={colback=iiiepeNavy,colframe=iiiepeNavy,arc=1pt},
  title={Reporte de lectura},#1}

\newtcolorbox{fichalecturabox}[1][]{%
  enhanced,breakable,
  colback=iiiepeAqua!7,colframe=iiiepeAqua!70!iiiepeNavy,
  boxrule=.6pt,arc=1pt,
  left=7pt,right=7pt,top=4pt,bottom=4pt,
  fontupper=\small,before upper={\setstretch{1.0}},
  fonttitle=\bfseries\color{iiiepeNavy},coltitle=iiiepeNavy,
  colbacktitle=iiiepeAqua!7,
  title={Ficha de la lectura},#1}



\usepackage{tabularx,colortbl,pgfplots,enumitem,needspace}\pgfplotsset{compat=1.18}\definecolor{iiiepePale}{HTML}{EEF5F8}
\setlength{\parskip}{5pt}\setlength{\parindent}{0pt}
\begin{document}\insertcoverwatermark\templatePortrait\templatePagecfg\templateFinalcfg\fancyfoot{}\fancyfoot[L]{\scriptsize IIIEPE · MAEM}\fancyfoot[C]{\scriptsize\thepage}\fancyfoot[R]{\scriptsize Actividad 11}\setstretch{1.3}\raggedbottom
\section{Problema y conexión matemática}
La sesión vincula sistemas de ecuaciones, perímetros y funciones mediante programación de un robot. El propósito es traducir relaciones algebraicas a instrucciones y contrastar un modelo con mediciones. En mantenimiento y mecatrónica, programar movimientos exige distinguir una solución teórica de la ejecución física, afectada por fricción, batería, superficie y calibración.
El recurso menciona soluciones de 14 y 22 segundos, pero no muestra en su texto un sistema completo. Se construye a continuación un ejemplo didáctico consistente con esos valores, identificado como propuesta propia y no como transcripción de un problema omitido.
\section{Reto 1 · Sistema de dos ecuaciones}
Sea $x$ el tiempo de avance y $y$ el tiempo de giro en dos programas. Se propone el sistema
\[
\begin{cases}x+y=36,\\2x+y=50.\end{cases}
\]
Al restar la primera ecuación de la segunda, $x=14$ segundos; por sustitución $y=36-14=22$ segundos. Las comprobaciones son $14+22=36$ y $2(14)+22=50$. La solución es única porque las ecuaciones son independientes. Los tiempos de giro deben calibrarse: 22 segundos no equivalen por sí mismos a un ángulo determinado.
\section{Reto 2 · Perímetro de una pista}
Para un trapecio isósceles de bases $b+4$ y $b$ y lados de 5 cm, $P=(b+4)+b+10=2b+14$. Si se requiere $P=30$ cm, $b=8$ y las bases miden 12 y 8 cm. La altura es $\sqrt{5^2-2^2}=\sqrt{21}$ cm; los giros exteriores deben calcularse a partir de la geometría y calibrarse físicamente. Un perímetro correcto no garantiza que cualquier combinación de giros cierre la trayectoria.
El algoritmo propuesto es: fijar origen y orientación, avanzar 12 cm, girar el ángulo exterior correspondiente, avanzar 5 cm, avanzar la base de 8 cm después del giro y cerrar con el segundo lado. Se verificaría posición final y error de cierre. Se dispone de pruebas físicas; este desarrollo explicita el algoritmo geométrico y se refiere cualitativamente a la relación entre programación y ejecución.
\section{Reto 3 · Calibración y función}
Con velocidad aproximadamente constante y origen temporal definido, $d(t)=vt$. La tabla es un ejemplo simulado que permite calcular $v=d/t=15$ cm/s; no representa un robot medido.


{\small\setstretch{1.1}\renewcommand{\arraystretch}{1.25}
\begin{longtable}{|>{\raggedright\arraybackslash}p{\dimexpr0.3\linewidth-2\tabcolsep-2\arrayrulewidth\relax}|>{\raggedright\arraybackslash}p{\dimexpr0.35\linewidth-2\tabcolsep-2\arrayrulewidth\relax}|>{\raggedright\arraybackslash}p{\dimexpr0.35\linewidth-2\tabcolsep-2\arrayrulewidth\relax}|}\hline
\rowcolor{iiiepePale} \textbf{Tiempo (s)} & \textbf{Distancia simulada (cm)} & \textbf{Cociente (cm/s)} \\\hline
\endfirsthead
\hline
\rowcolor{iiiepePale} \textbf{Tiempo (s)} & \textbf{Distancia simulada (cm)} & \textbf{Cociente (cm/s)} \\\hline
\endhead
2 & 30 & 15 \\\hline
4 & 60 & 15 \\\hline
6 & 90 & 15 \\\hline
\end{longtable}}


Si una medición arrojara $d=15t+4$, el término de 4 cm exigiría revisar origen, retardo o sesgo; ya no sería proporcionalidad directa. Para recorrer 120 cm con el modelo ideal $d=15t$, se requieren 8 s. Una prueba de transferencia puede cambiar la batería o la superficie y solicitar recalibración sin alterar los cálculos para ocultar el error.
\section{Completar la geometría y el algoritmo}
El trapecio propuesto tiene bases 12 y 8 cm, lados 5 cm y altura $\sqrt{21}$ cm. El ángulo agudo de la base satisface $\cos\alpha=2/5$, de modo que $\alpha\approx66.42^\circ$. Una trayectoria que recorra la base mayor, un lado, la base menor y el otro lado requiere giros exteriores de $113.58^\circ$, $66.42^\circ$, $66.42^\circ$ y $113.58^\circ$, cuya suma es $360^\circ$.
El algoritmo de referencia es: fijar origen y orientación; avanzar 12 cm; girar $113.58^\circ$; avanzar 5 cm; girar $66.42^\circ$; avanzar 8 cm; girar $66.42^\circ$; avanzar 5 cm; girar $113.58^\circ$. El sentido de giro se mantiene. La solución geométrica cierra una figura ideal; el robot necesita calibración independiente de avance y giro.
Para un avance con velocidad medida $v$, el tiempo es $t=d/v$. Para un giro aproximadamente uniforme de velocidad angular $\omega$, es $t_g=\theta/\omega$. No se utiliza el mismo parámetro para ambos. Se propone repetir pruebas de avance y giro, registrar error final de posición y orientación y revisar solo una condición por vez. No se presentan estos parámetros como resultados medidos.


\subsection{Referencia cualitativa a las pruebas del robot}
Las pruebas disponibles permiten situar la reflexión en la relación entre instrucciones y ejecución física. Su lectura se organiza en torno a la coherencia de la secuencia, la independencia de avance y giro y la justificación del ajuste. El modelo geométrico ofrece una referencia para analizar el comportamiento; la presencia de una trayectoria ejecutada no basta para atribuir precisión o mejora sin consultar los registros. En esta versión se conserva la referencia cualitativa solicitada y no se inventan tiempos, ángulos medidos ni errores de cierre.
El sistema numérico del desarrollo sigue siendo un ejemplo propio. La ausencia del sistema completo en el texto del recurso debe conservarse como límite de correspondencia, en lugar de atribuir al docente una ecuación que no se observó.
\section{Secuencia de aplicación y evaluación}
Vincular sistemas, geometría y calibración robótica. La duración propuesta es de 100 minutos, ajustable al diagnóstico y a los recursos. Esta planeación sigue la progresión de los recursos institucionales (IIIEPE, 2026); los tiempos y apoyos son decisiones de codiseño, no datos de una aplicación ya realizada.

{\small\setstretch{1.1}\renewcommand{\arraystretch}{1.25}
\begin{longtable}{|>{\raggedright\arraybackslash}p{\dimexpr0.23\linewidth-2\tabcolsep-2\arrayrulewidth\relax}|>{\raggedright\arraybackslash}p{\dimexpr0.49\linewidth-2\tabcolsep-2\arrayrulewidth\relax}|>{\raggedright\arraybackslash}p{\dimexpr0.28\linewidth-2\tabcolsep-2\arrayrulewidth\relax}|}\hline
\rowcolor{iiiepePale} \textbf{Fase y tiempo} & \textbf{Actividad y mediación} & \textbf{Evidencia} \\\hline
\endfirsthead
\hline
\rowcolor{iiiepePale} \textbf{Fase y tiempo} & \textbf{Actividad y mediación} & \textbf{Evidencia} \\\hline
\endhead
1. Reto y saberes (10 min) & Identificar qué sabe el modelo y qué requiere medir el robot. & Predicción individual y pregunta. \\\hline
2. Plan de indagación (20 min) & Resolver el sistema y deducir lados y giros de la pista. & Plan explicado y criterios. \\\hline
3. Construcción y análisis (35 min) & Programar avances y giros con calibraciones diferenciadas. & Registro o procedimiento trazable. \\\hline
4. Comunicación y aplicación (20 min) & Medir cierre, comparar predicción y ejecución y registrar el error. & Argumento con comprobación. \\\hline
5. Metacognición (15 min) & Revisar un parámetro y justificar el ajuste sin alterar datos. & Versión revisada y justificación. \\\hline
\end{longtable}}
\subsection{Criterios observables y retroalimentación}

{\small\setstretch{1.1}\renewcommand{\arraystretch}{1.25}
\begin{longtable}{|>{\raggedright\arraybackslash}p{\dimexpr0.22\linewidth-2\tabcolsep-2\arrayrulewidth\relax}|>{\raggedright\arraybackslash}p{\dimexpr0.4\linewidth-2\tabcolsep-2\arrayrulewidth\relax}|>{\raggedright\arraybackslash}p{\dimexpr0.38\linewidth-2\tabcolsep-2\arrayrulewidth\relax}|}\hline
\rowcolor{iiiepePale} \textbf{Criterio} & \textbf{Evidencia de logro} & \textbf{Acción si necesita apoyo} \\\hline
\endfirsthead
\hline
\rowcolor{iiiepePale} \textbf{Criterio} & \textbf{Evidencia de logro} & \textbf{Acción si necesita apoyo} \\\hline
\endhead
Sistema comprobado & Explica sistema comprobado con un caso de la actividad. & Pedir una representación concreta y contrastar dos casos. \\\hline
Geometría y giros & Conserva geometría y giros en el desarrollo y justifica sus pasos. & Modelar un paso, retirar gradualmente la ayuda y solicitar otro caso. \\\hline
Calibración documentada & Utiliza calibración documentada para contrastar una respuesta o representación. & Comparar la evidencia inicial y la corregida con una pregunta específica. \\\hline
Revisión del algoritmo & Relaciona revisión del algoritmo con la conclusión y un límite del problema. & Solicitar una comprobación independiente y un caso de transferencia. \\\hline
\end{longtable}}
El instrumento es una propuesta docente. Cada criterio se registra como logrado con autonomía, logrado con apoyo o aún no demostrado; no se asigna una calificación observada sin aplicar la actividad. En la coevaluación, un compañero verifica un criterio y formula una pregunta; en la autoevaluación, el estudiante explica una corrección; el docente contrasta ambas evidencias y decide el siguiente apoyo.\subsection{Participación y apoyos}
La inclusión se concreta en alternativas de representación y de acceso; el pensamiento crítico, en justificar y contrastar. Los roles de cálculo, registro y explicación se rotan para que todos produzcan una evidencia individual. Quien necesita apoyo dispone de un ejemplo parcial y preguntas graduadas; quien ya domina el procedimiento analiza un contraejemplo o cambia una condición. Estas acciones deben observarse durante la aplicación para valorar su funcionamiento.

\section{Secuencia y evaluación formativa}
En el inicio se identifican variables, puertos y criterios de éxito. En el desarrollo se resuelven los retos, se programan movimientos y se registra el error. En el cierre se compara predicción y medición. La autoevaluación explica una decisión; la coevaluación revisa código y cálculo; la heteroevaluación valora modelo, comprobación, unidades, participación e interpretación del error.
Los roles deben rotarse para evitar que la experiencia técnica de unos estudiantes concentre la programación. Puede utilizarse un simulador o un robot de papel con cuadrícula cuando no exista equipo suficiente. Esta alternativa conserva la argumentación matemática, aunque no permite evaluar calibración física.
\section{Reflexión pedagógica}
La robótica tiene sentido cuando obliga a revisar una relación matemática y no únicamente cuando produce movimiento. Mi trayectoria técnica facilita diseñar retos verificables; el acercamiento humanista me lleva a considerar la participación, el derecho a equivocarse y la explicación de cada estudiante. Propondría conservar una primera versión del cálculo y otra corregida, para valorar cómo cambió el razonamiento. La referencia a las pruebas reales es cualitativa: se examinan la separación entre parámetros de avance y giro, las condiciones de calibración y el criterio para revisar el cierre. Los valores del ejemplo no se atribuyen a esas pruebas.

\needspace{5\baselineskip}\section*{Referencias}
{\setstretch{1.15}\setlength{\parskip}{6pt}
\par\hangindent=1.27cm\hangafter=1  IIIEPE. (2026). \textit{Sesión 11: Ecuaciones, funciones y robótica}. Presentación proporcionada en EVA posgrado.
\par\hangindent=1.27cm\hangafter=1  Secretaría de Educación Pública. (2024). \textit{Programa de estudio para la educación secundaria: Programa sintético de la Fase 6}. SEP.
}

\end{document}
